\documentclass[10pt,letterpaper]{article}

\usepackage[margin=0.75in]{geometry}
\usepackage{amsmath,amssymb}
\usepackage{booktabs}
\usepackage{array}
\usepackage[dvipsnames]{xcolor}
\usepackage{enumitem}

\emergencystretch=1em

\newcommand{\op}[1]{\textsc{#1}}
\newcommand{\code}[1]{\texttt{\detokenize{#1}}}

\definecolor{callblue}{HTML}{1565C0}
\definecolor{assignbrown}{HTML}{795548}
\definecolor{markerorange}{HTML}{E65100}

\title{tinygrad: Lowering Pipeline}
\author{tinygrad, Corp. \\ \texttt{research@tinygrad.org}}
\date{}

\begin{document}
\maketitle
\thispagestyle{empty}

\section*{Overview}

tinygrad lowering is really two coupled pipelines:

\begin{enumerate}[leftmargin=*]
\item \textbf{Whole-graph lowering and scheduling}: turn a tensor graph into a scheduled list of executable kernels.
\item \textbf{Per-kernel lowering and code generation}: turn each scheduled kernel into \op{Program}, \op{Source}, and \op{Binary}.
\end{enumerate}

At a high level, the forms are:

\[
\text{tensor graph} \rightarrow \text{CALL graph} \rightarrow \text{kernel graph} \rightarrow \op{Linear} \rightarrow \text{ExecItem list}
\]

\[
\text{kernel } \op{Sink} \rightarrow \text{lowered kernel } \op{Sink} \rightarrow \op{Program}(\op{Linear}, \op{Source}, \op{Binary}) \rightarrow \text{CompiledRunner}
\]

\section*{Whole-Graph Lowering and Scheduling}

\begin{enumerate}[leftmargin=*]
\item \textbf{Build the root sink.} 
The entry points \code{Tensor.schedule_with_vars}, \code{Tensor.schedule}, and \code{Tensor.realize} start from a root \code{UOp.sink(...)} over the requested tensors.

\item \textbf{Callify the tensor graph.}
\code{transform_to_call} rewrites the tensor graph into a \op{Call} graph.
Its ordered passes are:
\begin{enumerate}[leftmargin=2em]
\item \code{number the uops}
\item \code{early transform tensor graph}
\item \code{finalize call}
\item \code{replace bufs}
\end{enumerate}
The result is a \op{Call} whose body is a sink of realized assignments with explicit \op{Param} inputs.

\item \textbf{Build the kernel graph.}
\code{complete_create_schedule_with_vars} rewrites each scheduled function sink through \code{pm_schedule}; this dispatches to \code{lower_sink_to_linear}, which first calls \code{get_kernel_graph}.

\code{get_kernel_graph} applies the following stages:
\begin{enumerate}[leftmargin=2em]
\item \code{multi_pm}
\item optional \code{fold moved afters}
\item \code{earliest rewrites}
\item \code{run_rangeify}
\item \code{symbolic+reduce_collapse+debuf}
\item \code{limit buffers}
\item \code{bufferize to store}
\item \code{split kernels}
\item optional \code{fix_assign}
\end{enumerate}

The output form is a kernel dependency graph built from \op{After}, \op{Call}, \op{Store}, and \op{End}. Each kernel is now an explicit schedulable unit.

\item \textbf{Rangeify the graph.}
The semantic center of whole-graph lowering is \code{run_rangeify}. It:
\begin{enumerate}[leftmargin=2em]
\item computes a realize map,
\item builds a consumer map,
\item propagates output ranges backward from consumers,
\item creates new \op{Range} axes when values must be realized,
\item turns \code{REDUCE_AXIS} into \op{Reduce} with explicit reduce ranges,
\item converts movement ops into index arithmetic, and
\item inserts \op{Bufferize} and \op{Index} where realized values cross kernel boundaries.
\end{enumerate}

\item \textbf{Linearize kernel dependencies.}
\code{create_schedule} topologically sorts the kernel dependency graph and returns a single \op{Linear} UOp. Each element of the \op{Linear} is a kernel \op{Call} in execution order.

\item \textbf{Resolve outer calls and allocate buffers.}
\code{pm_resolve_linear_call} resolves \code{CALL(LINEAR, ...)} by running \code{params to buffers}, allocating concrete \op{Buffer} objects, and flattening nested \op{Linear} nodes. Then \code{memory_plan_rewrite} performs memory planning across the linear schedule.

\item \textbf{Emit the executable schedule.}
\code{linear_to_schedule} converts the final \op{Linear} into a Python list of \code{ExecItem} objects. Each \code{ExecItem} contains the kernel AST, its buffers, metadata, and any fixed variable bindings.
\end{enumerate}

\section*{Per-Kernel Lowering and Code Generation}

\begin{enumerate}[leftmargin=*]
\item \textbf{Lower an ExecItem.}
\code{ExecItem.lower} dispatches compute kernels through \code{get_runner}. For \op{Sink}, \op{Program}, and \op{Beam} kernels, this calls \code{get_program} with the current renderer.

\item \textbf{Rewrite the kernel sink.}
If the input is a kernel \op{Sink}, \code{get_program} first calls \code{full_rewrite_to_sink}. This is the main per-kernel lowering pipeline.

Its ordered passes are:
\begin{enumerate}[leftmargin=2em]
\item \code{early movement ops}
\item \code{load collapse}
\item \code{split ranges}
\item \code{initial symbolic}
\item \code{simplify ranges}
\item \code{apply_opts}
\item \code{postopt symbolic}
\item \code{expander}
\item \code{add local buffers}
\item \code{remove_reduce}
\item \code{add gpudims}
\item \code{add loads (code)}
\item optional \code{create image buffers}
\item \code{devectorize}
\item \code{lower all index dtypes}
\item \code{post index symbolic}
\item optional \code{pre_matcher}
\item \code{decompositions}
\item \code{decomp dtypes}
\item \code{transcendental}
\item \code{final rewrite}
\item \code{add control flow}
\end{enumerate}

The output is still a kernel \op{Sink}, but now it is in renderer-facing program form: explicit loads, stores, control flow, specials, local buffers, and lowered arithmetic.

\item \textbf{Apply schedule-level kernel optimizations.}
The \code{apply_opts} stage runs the postrange scheduler. Depending on context, it can:
\begin{enumerate}[leftmargin=2em]
\item apply explicit \code{opts_to_apply},
\item run BEAM search,
\item or apply hand-coded heuristics.
\end{enumerate}
These transforms produce the final axis layout used by code generation: \texttt{GLOBAL}, \texttt{LOCAL}, \texttt{WARP}, \texttt{THREAD}, \texttt{REDUCE}, \texttt{GROUP\_REDUCE}, \texttt{UPCAST}, and \texttt{UNROLL}.

\item \textbf{Materialize \op{Program}.}
After kernel rewriting, \code{get_program} wraps the result as \code{PROGRAM(full\_sink, DEVICE)} and runs \code{pm_to_program}. This expands the program in four stages:
\begin{enumerate}[leftmargin=2em]
\item \code{do_linearize}: add \op{Linear}
\item \code{do_estimates}: compute FLOP and memory estimates
\item \code{do_render}: add \op{Source}
\item \code{do_compile}: add \op{Binary}
\end{enumerate}

The resulting form is:

\[
\op{Program}(\op{Sink}, \op{Device}, \op{Linear}, \op{Source}, \op{Binary})
\]

\item \textbf{Build the runtime program.}
\code{ProgramSpec.from_uop} extracts launch dimensions, runtime variables, globals, inputs, and outputs from the \op{Program}. \code{CompiledRunner} then owns the compiled program together with its launch metadata.

\item \textbf{Execute the schedule.}
\code{run_schedule} lowers each \code{ExecItem} on demand, ensures its buffers are allocated, launches the runner, and updates performance statistics.
\end{enumerate}

\section*{Canonical Stage Names}

The following names map cleanly to a spec because they denote stable transitions between graph forms:

\begin{enumerate}[leftmargin=*]
\item \code{transform_to_call}
\item \code{get_kernel_graph}
\item \code{run_rangeify}
\item \code{create_schedule}
\item \code{resolve linear call}
\item \code{memory plan}
\item \code{apply_opts}
\item \code{full_rewrite_to_sink}
\item \code{expander}
\item \code{devectorize}
\item \code{add gpudims}
\item \code{add control flow}
\item \code{linearize/render}
\item \code{ProgramSpec.from_uop}
\end{enumerate}

\section*{Pipeline Summary}

\begin{tabular}{@{}p{0.23\linewidth} p{0.27\linewidth} p{0.42\linewidth}@{}}
\toprule
\textbf{Stage} & \textbf{Input $\rightarrow$ Output} & \textbf{Main operation} \\
\midrule
Callify & tensor graph $\rightarrow$ \op{Call} graph & Realize assignments and replace external buffers with \op{Param}. \\
Kernel graph & \op{Call} graph $\rightarrow$ kernel graph & Rangeify, bufferize, split stores into schedulable kernels. \\
Schedule & kernel graph $\rightarrow$ \op{Linear} & Topologically order kernel dependencies. \\
Schedule materialization & \op{Linear} $\rightarrow$ \code{list[ExecItem]} & Allocate buffers, memory-plan, and attach runtime metadata. \\
Kernel rewrite & kernel \op{Sink} $\rightarrow$ lowered kernel \op{Sink} & Apply symbolic, layout, vectorization, reduction, and renderer-facing lowering passes. \\
Program materialization & lowered kernel \op{Sink} $\rightarrow$ \op{Program} & Linearize, estimate, render source, and compile binary. \\
Runtime & \op{Program} $\rightarrow$ executable runner & Extract launch metadata and execute on the target device. \\
\bottomrule
\end{tabular}

\end{document}
